% GENERATED FILE. Edit canonical lessons, then run npm run generate:books.
% canonical-source-hash: fnv1a64:7af81451c5312b8e
% canonical-lessons: GU-C443-jhad, GU-C443-vad, GU-C443-piplo, GU-C443-abo, GU-C443-phutvu

\chapter{Trees}
\label{ch:gu-trees}
\hlchaptermodality{443}

\begin{chapteropening}
\textbf{By the end of this chapter:} \emph{I can name the big trees of a Gujarati village.}

Complete the last lesson of chapter 443: I can name the big trees of a Gujarati village.
\end{chapteropening}

\section[\texorpdfstring{\gu{ઝાડ} (jhāḍ)}{ઝાડ (jhāḍ)}]{\gu{ઝાડ} (jhāḍ) — a tree}

\label{lesson:GU-C443-jhad}



\begin{quote}
Before the new one: say the Gujarati for a fishing net, then the Gujarati for to rock gently.
\end{quote}

\subsection*{You'll want to know: \gu{ઝાડ}}
\textbf{\gu{ઝાડ}} — \emph{jhāḍ} — \textquotedblleft{}a tree\textquotedblright{}.

\begin{quote}
\gu{અમારા} \gu{આંગણામાં} \gu{લીમડાનું} \gu{મોટું} \gu{ઝાડ} \gu{છે}.
\end{quote}

\begin{quote}
\emph{amārā ā̃gṇāmā̃ līmḍānũ moṭũ jhāḍ che.}
\end{quote}

\begin{quote}
There is a big neem tree in our courtyard.
\end{quote}

A neuter noun.

\begin{cousinweb}
\textbf{\gu{ઝાડ}} (\emph{jhāḍ}) is a tadbhava of Sanskrit \emph{jhāṭa}, \textquotedblleft{}a thicket, an arbour\textquotedblright{}; CDIAL lists the Gujarati word, \textquotedblleft{}a tree\textquotedblright{} (CDIAL 5362).
\end{cousinweb}

\begin{grammarlens}[title={What you've built}]
One more word for getting around.
\end{grammarlens}

\subsection*{Your turn}
\begin{itemize}
\raggedright
  \item \emph{Say it:} \emph{jhāḍ}
  \item \emph{Say it:} \emph{jhāḍ}, once more
  \item \emph{Say it:} \emph{amārā ā̃gṇāmā̃ līmḍānũ moṭũ jhāḍ che.} — There is a big neem tree in our courtyard.
  \item \emph{Recall:} say the Gujarati for a fishing net, then the Gujarati for to rock gently, then say \emph{jhāḍ} again
\end{itemize}

\begin{tcolorbox}[breakable,colback=teal!4,colframe=teal!35!black,title={Before you move on}]
What does \gu{ઝાડ} mean? (\textquotedblleft{}A tree\textquotedblright{}.) Say it once more.
\end{tcolorbox}
\section[\texorpdfstring{\gu{વડ} (vaḍ)}{વડ (vaḍ)}]{\gu{વડ} (vaḍ) — a banyan tree}

\label{lesson:GU-C443-vad}



\begin{quote}
Before the new one: say the Gujarati for to rock gently, then the Gujarati for a tree.
\end{quote}

\subsection*{You'll want to know: \gu{વડ}}
\textbf{\gu{વડ}} — \emph{vaḍ} — \textquotedblleft{}a banyan tree\textquotedblright{}.

\begin{quote}
\gu{ગામના} \gu{ચોકમાં} \gu{જૂનો} \gu{વડ} \gu{છે}.
\end{quote}

\begin{quote}
\emph{gāmnā cokmā̃ jūno vaḍ che.}
\end{quote}

\begin{quote}
There is an old banyan in the village square.
\end{quote}

A masculine noun.

\begin{cousinweb}
\textbf{\gu{વડ}} (\emph{vaḍ}) is a tadbhava of Sanskrit \emph{vaṭa}, \textquotedblleft{}the banyan\textquotedblright{}; CDIAL lists the Gujarati word (CDIAL 11211).
\end{cousinweb}

\begin{grammarlens}[title={What you've built}]
One more word for getting around.
\end{grammarlens}

\subsection*{Your turn}
\begin{itemize}
\raggedright
  \item \emph{Say it:} \emph{vaḍ}
  \item \emph{Say it:} \emph{vaḍ}, once more
  \item \emph{Say it:} \emph{gāmnā cokmā̃ jūno vaḍ che.} — There is an old banyan in the village square.
  \item \emph{Recall:} say the Gujarati for to rock gently, then the Gujarati for a tree, then say \emph{vaḍ} again
\end{itemize}

\begin{tcolorbox}[breakable,colback=teal!4,colframe=teal!35!black,title={Before you move on}]
What does \gu{વડ} mean? (\textquotedblleft{}A banyan tree\textquotedblright{}.) Say it once more.
\end{tcolorbox}
\section[\texorpdfstring{\gu{પીપળો} (pīpḷo)}{પીપળો (pīpḷo)}]{\gu{પીપળો} (pīpḷo) — a peepal tree}

\label{lesson:GU-C443-piplo}



\begin{quote}
Before the new one: say the Gujarati for a tree, then the Gujarati for a banyan.
\end{quote}

\subsection*{You'll want to know: \gu{પીપળો}}
\textbf{\gu{પીપળો}} — \emph{pīpḷo} — \textquotedblleft{}a peepal tree\textquotedblright{}.

\begin{quote}
\gu{મંદિર} \gu{પાસે} \gu{એક} \gu{પીપળો} \gu{છે}.
\end{quote}

\begin{quote}
\emph{mandir pāse ek pīpḷo che.}
\end{quote}

\begin{quote}
There is a peepal tree by the temple.
\end{quote}

A masculine noun.

\begin{cousinweb}
\textbf{\gu{પીપળો}} (\emph{pīpḷo}) is a tadbhava of Sanskrit \emph{pippala}, \textquotedblleft{}the sacred fig tree\textquotedblright{}; CDIAL lists the Gujarati word (CDIAL 8205).
\end{cousinweb}

\begin{grammarlens}[title={What you've built}]
One more word for getting around.
\end{grammarlens}

\subsection*{Your turn}
\begin{itemize}
\raggedright
  \item \emph{Say it:} \emph{pīpḷo}
  \item \emph{Say it:} \emph{pīpḷo}, once more
  \item \emph{Say it:} \emph{mandir pāse ek pīpḷo che.} — There is a peepal tree by the temple.
  \item \emph{Recall:} say the Gujarati for a tree, then the Gujarati for a banyan, then say \emph{pīpḷo} again
\end{itemize}

\begin{tcolorbox}[breakable,colback=teal!4,colframe=teal!35!black,title={Before you move on}]
What does \gu{પીપળો} mean? (\textquotedblleft{}A peepal tree\textquotedblright{}.) Say it once more.
\end{tcolorbox}
\section[\texorpdfstring{\gu{આંબો} (ā̃bo)}{આંબો (ā̃bo)}]{\gu{આંબો} (ā̃bo) — a mango tree}

\label{lesson:GU-C443-abo}



\begin{quote}
Before the new one: say the Gujarati for a banyan, then the Gujarati for a peepal.
\end{quote}

\subsection*{You'll want to know: \gu{આંબો}}
\textbf{\gu{આંબો}} — \emph{ā̃bo} — \textquotedblleft{}a mango tree\textquotedblright{}.

\begin{quote}
\gu{ઉનાળામાં} \gu{આંબા} \gu{પર} \gu{કેરીઓ} \gu{પાકે} \gu{છે}.
\end{quote}

\begin{quote}
\emph{unāḷāmā̃ ā̃bā par kerīo pāke che.}
\end{quote}

\begin{quote}
In summer mangoes ripen on the mango tree.
\end{quote}

A masculine noun.

\begin{cousinweb}
\textbf{\gu{આંબો}} (\emph{ā̃bo}) is a tadbhava of Sanskrit \emph{āmra}, \textquotedblleft{}the mango tree\textquotedblright{}; CDIAL lists the Gujarati word (CDIAL 1268). The fruit is \gu{કેરી}.
\end{cousinweb}

\begin{grammarlens}[title={What you've built}]
One more word for getting around.
\end{grammarlens}

\subsection*{Your turn}
\begin{itemize}
\raggedright
  \item \emph{Say it:} \emph{ā̃bo}
  \item \emph{Say it:} \emph{ā̃bo}, once more
  \item \emph{Say it:} \emph{unāḷāmā̃ ā̃bā par kerīo pāke che.} — In summer mangoes ripen on the mango tree.
  \item \emph{Recall:} say the Gujarati for a banyan, then the Gujarati for a peepal, then say \emph{ā̃bo} again
\end{itemize}

\begin{tcolorbox}[breakable,colback=teal!4,colframe=teal!35!black,title={Before you move on}]
What does \gu{આંબો} mean? (\textquotedblleft{}A mango tree\textquotedblright{}.) Say it once more.
\end{tcolorbox}
\section[\texorpdfstring{\gu{ફૂટવું} (phūṭvũ)}{ફૂટવું (phūṭvũ)}]{\gu{ફૂટવું} (phūṭvũ) — to sprout; to burst out}

\label{lesson:GU-C443-phutvu}



\begin{quote}
Before the new one: say the Gujarati for a peepal, then the Gujarati for a mango tree.
\end{quote}

\subsection*{You'll want to know: \gu{ફૂટવું}}
\textbf{\gu{ફૂટવું}} — \emph{phūṭvũ} — \textquotedblleft{}to sprout; to burst out\textquotedblright{}.

\begin{quote}
\gu{વસંતમાં} \gu{ઝાડ} \gu{પર} \gu{નવાં} \gu{પાન} \gu{ફૂટે} \gu{છે}.
\end{quote}

\begin{quote}
\emph{vasantmā̃ jhāḍ par navā̃ pān phūṭe che.}
\end{quote}

\begin{quote}
In spring new leaves sprout on the trees.
\end{quote}

Past (it …): \gu{તે} \gu{પાન} \gu{ફૂટ્યું}, \emph{te pān phūṭyũ}.

\begin{cousinweb}
\textbf{\gu{ફૂટવું}} (\emph{phūṭvũ}) goes back to \emph{*sphuṭyati}, \textquotedblleft{}splits open\textquotedblright{}, a form CDIAL reconstructs; CDIAL lists the Gujarati verb, \textquotedblleft{}to split, to burst\textquotedblright{} (CDIAL 13845): a bud that bursts into leaf.
\end{cousinweb}

\begin{grammarlens}[title={What you've built}]
One more word for getting around.
\end{grammarlens}

\subsection*{Your turn}
\begin{itemize}
\raggedright
  \item \emph{Say it:} \emph{phūṭvũ}
  \item \emph{Say it:} \emph{phūṭvũ}, once more
  \item \emph{Say it:} \emph{vasantmā̃ jhāḍ par navā̃ pān phūṭe che.} — In spring new leaves sprout on the trees.
  \item \emph{Recall:} say the Gujarati for a peepal, then the Gujarati for a mango tree, then say \emph{phūṭvũ} again
\end{itemize}

\begin{tcolorbox}[breakable,colback=teal!4,colframe=teal!35!black,title={Before you move on}]
What does \gu{ફૂટવું} mean? (\textquotedblleft{}To sprout; to burst out\textquotedblright{}.) Say it once more.
\end{tcolorbox}
